\SourceLeaf{126}%
$ay-bx=\eta$, erit $\eta\equiv k\xi$, $\xi\equiv-k\eta(\mathrm{mod.}\ p)$. Omnes itaque numeri $\xi+\eta i$, quibus residua simpliciter minima $x+yi$ respondent, habebuntur, dum vel pro $\xi$ deinceps accipiuntur valores $0,1,2,3\dotsfour p-1$, et pro $\eta$ residua minima positiva productorum $k\xi$ secundum modulum $p$, vel ordine alio pro $\eta$ illi valores et pro $\xi$ residua minima productorum $-k\eta$. E singulis $\xi+\eta i$ dein respondentes $x+yi$ invenientur per formulam
\[
x+yi=\frac{\xi+\eta i}{a-bi}=\frac{a\xi-b\eta}{p}+\frac{a\eta+b\xi}{p}\pdot i
\]

Ceterum obvium est, $\eta$, dum $\xi$ unitate crescat, vel augmentum $k$ vel decrementum $p-k$ pati, adeoque $x+yi$
\[
\begin{aligned}
&\text{vel mutationem }\frac{a-kb}{p}+\frac{ak+b}{p}\pdot i\\
&\text{vel hanc }\frac{a-kb}{p}+b+\left(\frac{ak+b}{p}-a\right)i
\end{aligned}
\]
quae observatio ad constructionem faciliorem reddendam inservit.

Denique patet, si residua absolute minima $x+yi$ desiderentur, haec praecepta eatenus tantum mutari, quatenus ipsi $\xi$ deinceps tribuendi sint valores inter limites $-\frac12p$ et $+\frac12p$, dum pro $\eta$ accipere oporteat residua absolute minima productorum $k\xi$. Ecce conspectum residuorum minimorum pro modulo $5+2i$ hoc modo adornatorum:

\begin{center}Residua simpliciter minima.\end{center}
{\small\setlength{\arraycolsep}{3pt}
\[
\begin{array}{r|r||r|r||r|r}
\xi+\eta i&x+yi&\xi+\eta i&x+yi&\xi+\eta i&x+yi\\\hline
0&0&10+25i&+5i&20+21i&+2+5i\\
1+17i&-1+3i&11+13i&+1+3i&21+9i&+3+3i\\
2+5i&+i&12+i&+2+i&22+26i&+2+6i\\
3+22i&+1+4i&13+18i&+1+4i&23+14i&+3+4i\\
4+10i&+2i&14+6i&+2+2i&24+2i&+4+2i\\
5+27i&-1+5i&15+23i&+1+5i&25+19i&+3+5i\\
6+15i&+3i&16+11i&+2+3i&26+7i&+4+3i\\
7+3i&+1+i&17+28i&+1+6i&27+24i&+3+6i\\
8+20i&+4i&18+16i&+2+4i&28+12i&+4+4i\\
9+8i&+1+2i&19+4i&+3+2i&&
\end{array}
\]
}

\SourceLeaf{127}%
\begin{center}Residua absolute minima.\end{center}
{\small\setlength{\arraycolsep}{3pt}
\[
\begin{array}{r|r||r|r||r|r}
\xi+\eta i&x+yi&\xi+\eta i&x+yi&\xi+\eta i&x+yi\\\hline
-14-6i&-2-2i&-4-10i&-2i&+5-2i&+1\\
-13+11i&-3+i&-3+7i&-1+i&+6-14i&+2-2i\\
-12-i&-2-i&-2-5i&-i&+7+3i&+1+i\\
-11-13i&-1-3i&-1+12i&-1+2i&+8-9i&+2-i\\
-10+4i&-2&0&0&+9+8i&+1+2i\\
-9-8i&-1-2i&+1-12i&+1-2i&+10-4i&+2\\
-8+9i&-2+i&+2+5i&+i&+11+13i&+1+3i\\
-7-3i&-1-i&+3-7i&+1-i&+12+i&+2+i\\
-6+14i&-2+2i&+4+10i&+2i&+13-11i&+3-i\\
-5+2i&-1&&&+14+6i&+2+2i
\end{array}
\]
}

Casum secundum, ubi $a,b$ non sunt inter se primi, facile ad casum praecedentem reducere licet. Sit $\lambda$ divisor communis maximus numerorum $a,b$, atque $a=\lambda a'$, $b=\lambda b'$. Denotet $F$ indefinite residuum minimum pro modulo $\lambda$, quatenus tamquam numerus complexus consideratur, i. e. exhibeat indefinite numerum talem $x+yi$, ut $x,y$ sint vel inter limites $0$ et $\lambda$, vel inter hos $-\frac12\lambda$ et $+\frac12\lambda$ (prout de residuis vel simpliciter vel absolute minimis agitur): denotet porro $F'$ indefinite residuum minimum pro modulo $a'+b'i$. Tunc erit $(a'+b'i)F+F'$ indefinite residuum minimum pro modulo $a+bi$, prodibitque systema completum horum residuorum, dum omnia $F$ cum omnibus $F'$ combinantur.

\Article{46}
Duo numeri complexi inter se primi dicuntur, si praeter unitates alios divisores communes non admittunt: quoties autem tales divisores communes adsunt, ii divisores communes maximi vocantur, quorum norma maxima est.

Si duorum numerorum propositorum resolutio in factores primos praesto est, determinatio divisoris communis maximi prorsus eodem modo perficitur, ut pro numeris realibus (\emph{Disquiss. Ar.} art. 18). Simul hinc elucet, omnes divisores communes duorum numerorum datorum metiri debere eorundem divisorem communem maximum hoc modo inventum. Quare quum sponte iam pateat, ternos numeros huic socios etiam esse divisores communes, semper quaterni numeri, et non plu%
\SourceLeaf{128}%
res, divisores communes maximi appellandi erunt, horumque norma erit multiplum normae cuiusvis alius divisoris communis.

Si resolutio duorum numerorum propositorum in factores simplices non adest, divisor communis maximus adiumento similis algorithmi eruitur, ut pro numeris realibus. Sint $m,m'$ duo numeri propositi, formeturque per divisionem repetitam series $m'',m'''$ etc. ita, ut $m''$ sit residuum absolute minimum ipsius $m$ secundum modulum $m'$, dein $m'''$ residuum absolute minimum ipsius $m'$ secundum modulum $m''$ et sic porro. Denotando normas numerorum $m,m',m'',m'''$ etc. resp. per $p,p',p'',p'''$ etc., erit $\frac{p''}{p'}$ norma quotientis $\frac{m''}{m'}$, adeoque per definitionem residui absolute minimi certo non maior quam $\frac12$; idem valet de $\frac{p'''}{p''}$ etc. Quapropter integri reales positivi $p',p'',p'''$ etc. seriem continuo decrescentem formabunt, unde necessario tandem ad terminum $0$ pervenietur, sive, quod idem est, in serie $m,m',m'',m'''$ etc. tandem ad terminum perveniemus, qui praecedentem absque residuo metitur. Sit hic $m^{(n+1)}$, statuamusque
\[
\begin{aligned}
m&=km'+m''\\
m'&=k'm''+m'''\\
m''&=k''m'''+m''''
\end{aligned}
\]
etc. usque ad
\[
m^{(n)}=k^{(n)}m^{(n+1)}
\]

Percurrendo has aequationes ordine inverso, elucet, $m^{(n+1)}$ singulos terminos praecedentes $m^{(n)}\dotsfour m'',m',m$ metiri; percurrendo autem easdem aequationes ordine directo, manifestum est, quemvis divisorem communem numerorum $m,m'$ etiam metiri singulos sequentes. Conclusio prior docet, $m^{(n+1)}$ esse divisorem communem numerorum $m,m'$; posterior autem, hunc divisorem esse maximum.

Ceterum quoties residuum ultimum $m^{(n+1)}$ alicui quatuor unitatum $1,-1,i,-i$ aequale evadit, hoc indicium erit, $m$ et $m'$ inter se primos esse.

\Article{47}
Si aequationes art. praec., omissa ultima, ita combinantur, ut $m'',\allowbreak m''',\allowbreak m''''\dotsfour\allowbreak m^{(n)}$ eliminentur, orietur aequatio talis
\[
m^{(n+1)}=hm+h'm'
\]
\SourceLeaf{129}%
ubi $h,h'$ erunt integri, et quidem, si designatione in \emph{Disquiss. Ar.} art. 27 introducta uti placet
\[
\begin{aligned}
h&=\pm[k',k'',k'''\dotsfour k^{(n-1)}]\\
&=\pm[k^{(n-1)},k^{(n-2)}\dotsfour k'',k']\\
h'&=\mp[k,k',k'',k'''\dotsfour k^{(n-1)}]\\
&=\mp[k^{(n-1)},k^{(n-2)}\dotsfour k'',k',k]
\end{aligned}
\]
valentibus signis superioribus vel inferioribus, prout $n$ par est vel impar. Hoc theorema ita enunciamus:

\emph{Divisor communis maximus duorum numerorum complexorum $m,m'$ redigi potest ad formam $hm+h'm'$, ita ut $h,h'$ sint integri.}

Manifesto enim hoc non solum de eo divisore communi maximo valet, ad quem algorithmus art. praec. deduxit, sed etiam de tribus illi associatis, pro quibus loco coëfficientium $h,h'$ accipere oportebit vel hos $hi,h'i$ vel $-h,-h'$, vel $-hi,-h'i$.

Quoties itaque numeri $m,m'$ inter se primi sunt, satisfieri poterit aequationi
\[
1=hm+h'm'
\]

Propositi sint e. g. numeri $31+6i=m$, $11-20i=m'$. Hic invenimus
\[
\begin{aligned}
k&=i,&m''&=+11-5i\\
k'&=+1-i,&m'''&=+5-4i\\
k''&=+2,&m''''&=+1+3i\\
k'''&=-1-2i,&m'''''&=+i\\
k''''&=+3-i
\end{aligned}
\]
atque hinc
\[
\begin{aligned}
[k',k'',k''']&=-6-5i\\
[k,k',k'',k''']&=+4-10i
\end{aligned}
\]
et proin
\[
m'''''=i=(6+5i)m+(4-10i)m'
\]
nec non
\[
1=(5-6i)m+(-10-4i)m'
\]
quod calculo instituto confirmatur.

\Article{48}
Per praecedentia omnia, quae ad theoriam congruentiarum primi gradus in arithmetica numerorum complexorum requiruntur, praeparata sunt: sed quum illa
\SourceLeaf{130}%
essentialiter non differat ab ea, quae pro arithmetica numerorum realium locum habet, atque in \emph{Disquisitionibus Arithmeticis} copiose exposita est, praecipua momenta hic adscripsisse sufficiet.

I. Congruentia $mt\equiv1\ (\mathrm{mod.}\ m')$ aequivalet aequationi indeterminatae $mt+m'u=1$, et si huic satisfit per valores $t=h,u=h'$, illius solutio generaliter exhibetur per $t\equiv h(\mathrm{mod.}\ m')$: conditio autem solubilitatis est, ut modulus $m'$ cum coëfficiente $m$ divisorem communem non habeat.

II. Solutio congruentiae $ax+b\equiv c(\mathrm{mod.}\ M)$ in casu eo, ubi $a,M$ sunt inter se primi, pendet a solutione huius
\[
at\equiv1\ (\mathrm{mod.}\ M)
\]
cui si satisfacit $t=h$, illius solutio generalis continetur in formula
\[
x\equiv(c-b)h\ (\mathrm{mod.}\ M)
\]

III, Congruentia $ax+b\equiv c(\mathrm{mod.}\ M)$ in casu eo, ubi $a,M$ divisorem communem $\lambda$ habent, aequivalet huic
\[
\frac a\lambda\pdot x\equiv\frac{c-b}\lambda\ \left(\mathrm{mod.}\frac M\lambda\right)
\]
Dum itaque pro $\lambda$ adoptatur divisor communis maximus numerorum $a,M$, solutio congruentiae propositae ad casum praecedentem reducitur, patetque, ad resolubilitatem requiri et sufficere, ut $\lambda$ etiam differentiam $c-b$ metiatur.

\Article{49}
Hactenus elementaria tantum attigimus, quae tamen nexus caussa omittere non licuit. In disquisitionibus altioribus arithmetica numerorum complexorum arithmeticae realium in eo similis est, quod theoremata elegantiora et simpliciora prodeunt, dum tales modulos, qui sunt numeri primi, solos admittimus: revera illorum extensio ad modulos compositos plerumque prolixior quam difficilior est, et laboris potius quam artis. Quapropter in sequentibus imprimis de modulis primis agetur.

\Article{50}
Denotante $X$ functionem indeterminatae $x$ talem
\[
Ax^n+Bx^{n-1}+Cx^{n-2}+\text{etc.}+Mx+N;
\]
